\section{Klasifikavimo pavyzdžiai}
\label{sec:pavyzdziu_klasifikacija}

Iš testavimo aibės parinkta po $30$ įrašų abiem užduotims; jiems geriausias modelis (žr.~\ref{sec:geriausias_modelis} sk.) suformulavo spėjimą. Įrašai paimti iš testavimo \texttt{DataLoader}'io pirmojo paketo (\texttt{batch\_size}~$=64$). Žaliai pažymėti teisingai klasifikuoti įrašai, raudonai – klaidingai.

\subsection{Vaizdų pavyzdžiai}
\label{subsec:pavyzdziai_vaizdai}

\begin{table}[H]
    \centering
    \caption{Vaizdų geriausio modelio $30$ pavyzdinių spėjimų testavimo aibėje.}
    \begin{tabular}{|c|l|l|c|}
        \hline
        Nr. & Tikra klasė & Spėjimas & Rezultatas \\
        \hline
        1  & muffin    & muffin    & \cellcolor{green!20}\checkmark \\ \hline
        2  & chihuahua & chihuahua & \cellcolor{green!20}\checkmark \\ \hline
        3  & chihuahua & chihuahua & \cellcolor{green!20}\checkmark \\ \hline
        4  & chihuahua & muffin    & \cellcolor{red!20}\texttimes \\ \hline
        5  & muffin    & muffin    & \cellcolor{green!20}\checkmark \\ \hline
        6  & chihuahua & chihuahua & \cellcolor{green!20}\checkmark \\ \hline
        7  & chihuahua & chihuahua & \cellcolor{green!20}\checkmark \\ \hline
        8  & chihuahua & chihuahua & \cellcolor{green!20}\checkmark \\ \hline
        9  & muffin    & chihuahua & \cellcolor{red!20}\texttimes \\ \hline
        10 & muffin    & muffin    & \cellcolor{green!20}\checkmark \\ \hline
        11 & chihuahua & chihuahua & \cellcolor{green!20}\checkmark \\ \hline
        12 & muffin    & muffin    & \cellcolor{green!20}\checkmark \\ \hline
        13 & muffin    & muffin    & \cellcolor{green!20}\checkmark \\ \hline
        14 & muffin    & muffin    & \cellcolor{green!20}\checkmark \\ \hline
        15 & muffin    & muffin    & \cellcolor{green!20}\checkmark \\ \hline
        16 & muffin    & muffin    & \cellcolor{green!20}\checkmark \\ \hline
        17 & muffin    & muffin    & \cellcolor{green!20}\checkmark \\ \hline
        18 & muffin    & muffin    & \cellcolor{green!20}\checkmark \\ \hline
        19 & chihuahua & chihuahua & \cellcolor{green!20}\checkmark \\ \hline
        20 & chihuahua & chihuahua & \cellcolor{green!20}\checkmark \\ \hline
        21 & chihuahua & chihuahua & \cellcolor{green!20}\checkmark \\ \hline
        22 & chihuahua & muffin    & \cellcolor{red!20}\texttimes \\ \hline
        23 & chihuahua & chihuahua & \cellcolor{green!20}\checkmark \\ \hline
        24 & muffin    & muffin    & \cellcolor{green!20}\checkmark \\ \hline
        25 & chihuahua & chihuahua & \cellcolor{green!20}\checkmark \\ \hline
        26 & chihuahua & chihuahua & \cellcolor{green!20}\checkmark \\ \hline
        27 & muffin    & muffin    & \cellcolor{green!20}\checkmark \\ \hline
        28 & muffin    & muffin    & \cellcolor{green!20}\checkmark \\ \hline
        29 & muffin    & muffin    & \cellcolor{green!20}\checkmark \\ \hline
        30 & chihuahua & chihuahua & \cellcolor{green!20}\checkmark \\ \hline
    \end{tabular}
    \label{tab:image_samples}
\end{table}


Iš $30$ pavyzdžių $27$ klasifikuoti teisingai ($90\,\%$), $3$ – klaidingai. Visos klaidos buvo \textit{muffin}~$\rightarrow$ \textit{chihuahua} arba atvirkščiai – tai nestebina, nes tai yra binarinė užduotis su dažnai vizualiai panašiais objektais.

\subsection{Laiko eilučių pavyzdžiai}
\label{subsec:pavyzdziai_keystroke}

\begin{table}[H]
    \centering
    \caption{Laiko eilučių geriausio modelio $30$ pavyzdinių spėjimų testavimo aibėje.}
    \begin{tabular}{|c|l|l|c|}
        \hline
        Nr. & Tikra klasė & Spėjimas & Rezultatas \\
        \hline
        1  & s008 & s012 & \cellcolor{red!20}\texttimes \\ \hline
        2  & s017 & s017 & \cellcolor{green!20}\checkmark \\ \hline
        3  & s029 & s029 & \cellcolor{green!20}\checkmark \\ \hline
        4  & s029 & s029 & \cellcolor{green!20}\checkmark \\ \hline
        5  & s031 & s031 & \cellcolor{green!20}\checkmark \\ \hline
        6  & s015 & s015 & \cellcolor{green!20}\checkmark \\ \hline
        7  & s008 & s020 & \cellcolor{red!20}\texttimes \\ \hline
        8  & s010 & s010 & \cellcolor{green!20}\checkmark \\ \hline
        9  & s003 & s003 & \cellcolor{green!20}\checkmark \\ \hline
        10 & s030 & s030 & \cellcolor{green!20}\checkmark \\ \hline
        11 & s007 & s007 & \cellcolor{green!20}\checkmark \\ \hline
        12 & s026 & s026 & \cellcolor{green!20}\checkmark \\ \hline
        13 & s026 & s026 & \cellcolor{green!20}\checkmark \\ \hline
        14 & s021 & s027 & \cellcolor{red!20}\texttimes \\ \hline
        15 & s026 & s026 & \cellcolor{green!20}\checkmark \\ \hline
        16 & s024 & s024 & \cellcolor{green!20}\checkmark \\ \hline
        17 & s029 & s029 & \cellcolor{green!20}\checkmark \\ \hline
        18 & s013 & s013 & \cellcolor{green!20}\checkmark \\ \hline
        19 & s025 & s025 & \cellcolor{green!20}\checkmark \\ \hline
        20 & s010 & s010 & \cellcolor{green!20}\checkmark \\ \hline
        21 & s032 & s032 & \cellcolor{green!20}\checkmark \\ \hline
        22 & s032 & s032 & \cellcolor{green!20}\checkmark \\ \hline
        23 & s015 & s015 & \cellcolor{green!20}\checkmark \\ \hline
        24 & s004 & s004 & \cellcolor{green!20}\checkmark \\ \hline
        25 & s011 & s011 & \cellcolor{green!20}\checkmark \\ \hline
        26 & s022 & s022 & \cellcolor{green!20}\checkmark \\ \hline
        27 & s002 & s002 & \cellcolor{green!20}\checkmark \\ \hline
        28 & s030 & s030 & \cellcolor{green!20}\checkmark \\ \hline
        29 & s020 & s020 & \cellcolor{green!20}\checkmark \\ \hline
        30 & s033 & s033 & \cellcolor{green!20}\checkmark \\ \hline
    \end{tabular}
    \label{tab:keystroke_samples}
\end{table}


Iš $30$ pavyzdžių $27$ klasifikuoti teisingai ($90\,\%$). Pavyzdžiai apima $21$ skirtingą klasę iš $30$ – ne visos klasės pateko į pirmąjį $30$ įrašų poaibį, nes \texttt{Trainer.get\_sample\_predictions} ima iš pirmųjų testavimo paketų, neužtikrindamas klasių subalansavimo.
